% Generated draft: requires mathematical and source review.
\documentclass[11pt]{article}
\usepackage[margin=25mm]{geometry}
\usepackage{fontspec}
\setmainfont{Latin Modern Roman}
\usepackage{amsmath,amssymb,mathtools,unicode-math}
\setmathfont{Latin Modern Math}
\usepackage{enumitem,needspace,xurl,hyperref}
\setlength{\parindent}{0pt}
\setlength{\parskip}{5pt}
\setlength{\emergencystretch}{3em}
\newcommand{\N}{\mathbb{N}}
\newcommand{\Z}{\mathbb{Z}}
\newcommand{\Q}{\mathbb{Q}}
\newcommand{\R}{\mathbb{R}}
\newcommand{\C}{\mathbb{C}}
\newcommand{\F}{\mathbb{F}}
\newcommand{\A}{\mathbb{A}}
\newcommand{\PP}{\mathbb{P}}
\newcommand{\E}{\mathbb{E}}
\newcommand{\eps}{\varepsilon}
\newcommand{\dd}{\,\mathrm{d}}
\begin{document}
\section*{UnsolvedMath ID: 30001299}
\textbf{Source problem code:} catalogue entry
\subsection{Definitions and mathematical statement}
Let \(k\) be a number field, that is, a finite extension of \(\Q\). For a variety \(X\) over \(k\), write
\[
X(k)=\operatorname{Hom}_k(\operatorname{Spec}k,X)
\]
for its set of \(k\)-rational points. Assume \(X(k)\ne\varnothing\). Regard these points as points of the underlying Zariski topological space of \(X\), and let \(\overline{X(k)}^{\mathrm{Zar}}\) denote their Zariski closure. The set \(X(k)\) is Zariski dense in \(X\) if
\[
\overline{X(k)}^{\mathrm{Zar}}=X.
\]
Equivalently, every proper Zariski closed subset of \(X\) omits at least one point of \(X(k)\).

\textbf{Classification question.} For every number field \(k\) and every variety \(X/k\) in the source's intended class of ``nice'' or ``suitably regular'' varieties with \(X(k)\ne\varnothing\), characterize the behavior of \(X(k)\). In particular, determine conditions under which
\[
1\leq \#X(k)<\infty,
\qquad
\#X(k)=\infty,
\]
and determine when \(X(k)\) is Zariski dense in \(X\). Density is a separate property, rather than a third alternative to finiteness and infinitude.

\textbf{Unresolved scope.} The supplied record does not define ``nice'' or ``suitably regular,'' fix conventions for ``variety,'' or specify further hypotheses such as smoothness, projectivity, or geometric integrality. Consequently, the intended class over which this question ranges remains unspecified. The source asks a broad classification question; it supplies neither a particular criterion to prove nor a precise additional meaning of ``distribution'' beyond the examples above.
\subsection{Short English statement}
For a variety over a number field that has at least one rational point, seek criteria describing whether its rational points are finite or infinite and whether they are Zariski dense. The intended regularity assumptions remain unspecified, so this is a classification question requiring clarification rather than a fully specified conjecture.
\subsection{Sources}
\begin{itemize}
\item Supplied local UnsolvedMath registry record, permanent ID \(30001299\), title ``Distribution of Rational Points on Varieties,'' source problem code \texttt{catalogue entry}, in the Oberwolfach Reports Open Problems collection. No individual proposer is supplied; the proposed year is recorded as 2009.
\item Supplied report citation: \emph{Explicit Methods in Number Theory (2010)}, Oberwolfach Reports, pp.\ 1843--1920, DOI \texttt{10.4171/owr/2009/33}. \url{https://doi.org/10.4171/owr/2009/33}. The record describes Oberwolfach Reports as documenting mathematical workshops and their open-problem discussions.
\item Supplied literature source labeled ``Poonen Oberwolfach report'': \url{https://ems.press/content/serial-article-files/46233?nt=1}.
\item Supplied literature source labeled ``Poonen survey'': \url{https://math.mit.edu/~poonen/papers/millennial.pdf}.
\item Literature assessment dated \textbf{2026-08-21}: the supplied curation labels the question open and describes it as a broad classification problem, with general behavior known only in important special classes and unresolved in general. This is supplied dated curation derived from the OWR record, not an independently verified literature assessment. Its time-sensitive status and cited claims require independent verification.
\item The supplied statement status is \texttt{corrected\_verified}; its accompanying verification note describes meaning-preserving semantic, typographic, and LaTeX cleanup without unsupported source-dependent additions. These are supplied registry qualifications, not independent verification of the formulation or mathematical claims.
\end{itemize}
\end{document}
